\documentclass{article}
\usepackage[preprint]{neurips_2024}
\usepackage[utf8]{inputenc}
\usepackage[T1]{fontenc}
\usepackage{hyperref}
\usepackage{url}
\usepackage{booktabs}
\usepackage{amsfonts}
\usepackage{amsmath}
\usepackage{amssymb}
\usepackage{amsthm}
\usepackage{nicefrac}
\usepackage{microtype}
\usepackage{graphicx}
\usepackage{xcolor}
\usepackage{algorithm}
\usepackage{algorithmic}
\usepackage{multirow}
\usepackage{subcaption}
\usepackage{mathtools}

\newtheorem{definition}{Definition}
\newtheorem{theorem}{Theorem}
\newtheorem{proposition}{Proposition}
\newtheorem{corollary}{Corollary}
\newtheorem{lemma}{Lemma}
\newtheorem{remark}{Remark}

\DeclareMathOperator*{\argmin}{arg\,min}
\DeclareMathOperator*{\argmax}{arg\,max}
\newcommand{\piref}{\pi_{\text{ref}}}
\newcommand{\pistar}{\pi^{*}}
\newcommand{\RR}{\mathbb{R}}
\newcommand{\EE}{\mathbb{E}}
\newcommand{\KL}{\mathrm{KL}}
\newcommand{\NashDPO}{\textsc{Nash-DPO}}

\title{Nash-DPO: Multi-Objective Alignment via Nash Bargaining\\with Game-Theoretic Self-Play}

\author{
  Anonymous Authors
}

\begin{document}
\maketitle

\begin{abstract}
Aligning large language models (LLMs) with human preferences involves balancing multiple competing objectives---correctness, safety, efficiency, and creativity---yet standard DPO and RLHF methods treat alignment as a single-objective problem.
We formalize multi-objective alignment as a \emph{Nash bargaining game} among $K$ objective agents and derive \NashDPO{}, a principled extension of DPO that adaptively weights objectives via the Nash bargaining solution.
We prove that \NashDPO{} converges to a Pareto-optimal policy satisfying individual rationality and symmetry axioms (Theorem~\ref{thm:pareto}), and that the optimal weights admit a closed-form KKT characterization that can be estimated online with $O(K)$ overhead (Theorem~\ref{thm:kkt}).
To enhance the strategic reasoning capabilities that underpin multi-agent alignment, we introduce a complementary GRPO self-play stage across four representative strategic games.
Through a full factorial experimental design (Base, SFT, A-only, B-only, A+B) on Qwen3.5-9B, we demonstrate that:
(1)~\NashDPO{} Pareto-dominates uniform-weight and fixed-weight DPO on 3 of 4 objectives while maximizing the Nash social welfare product;
(2)~GRPO self-play improves Nash equilibrium convergence by 23\% across games;
(3)~factorial ANOVA confirms a significant interaction effect ($p < 0.01$), with the combined A+B configuration achieving the best overall performance.
\end{abstract}


%====================================
\section{Introduction}
\label{sec:intro}
%====================================

Current large language model (LLM) post-training reduces human preferences to a single scalar reward, either through RLHF~\citep{ouyang2022training} or DPO~\citep{rafailov2024direct}.
This conflation obscures fundamental tensions: a response that maximizes helpfulness may compromise safety; one that maximizes conciseness may sacrifice creativity.
In practice, alignment engineers hand-tune reward weights or employ ad hoc multi-objective heuristics, with no guarantee of optimality~\citep{rame2024rewarded,wu2024fine}.

A second limitation is that standard post-training does not develop \emph{strategic reasoning}: the capacity to model other agents' beliefs and incentives and respond accordingly~\citep{gandhi2024strategic}.
Yet strategic reasoning is essential for real-world deployment where LLMs interact with other agents---human or artificial---in settings with competing interests~\citep{duan2024gtbench}.

We address both limitations through a game-theoretic lens.
Multi-objective alignment is naturally a \emph{cooperative bargaining game}: $K$ objective agents (correctness, safety, efficiency, creativity) must agree on a single policy $\pi$ such that each agent's utility exceeds a disagreement baseline.
The \emph{Nash bargaining solution}~\citep{nash1950bargaining,nash1953two} is the unique outcome satisfying Pareto optimality, individual rationality, independence of irrelevant alternatives, and symmetry.

Our contributions:
\begin{enumerate}
    \item \textbf{Formal framework.}
    We define the \emph{Multi-Objective Alignment Game} (MOAG, Definition~\ref{def:moag}) and prove that the Nash bargaining solution yields a Pareto-optimal policy with closed-form DPO weights (Theorems~\ref{thm:pareto}--\ref{thm:kkt}).
    
    \item \textbf{\NashDPO{} algorithm.}
    We translate the theoretical solution into a practical training algorithm (Algorithm~\ref{alg:nash_dpo}) with $O(K)$ overhead per step and EMA-stabilized weight updates.
    
    \item \textbf{GRPO self-play.}
    We show that self-play across representative strategic games develops transferable reasoning capabilities that complement multi-objective alignment.
    
    \item \textbf{Factorial validation.}
    A 5-condition factorial design with proper equal-compute baselines rigorously isolates each component's contribution and confirms significant complementarity via ANOVA interaction tests.
\end{enumerate}


%====================================
\section{Related Work}
\label{sec:related}
%====================================

\paragraph{Multi-objective alignment.}
Reward Soups~\citep{rame2024rewarded} interpolate between single-reward fine-tuned models in weight space.
MORLHF~\citep{wu2024fine} applies multi-objective RL with linear scalarization.
Both require choosing fixed weights; neither provides optimality guarantees.
Nash-MD~\citep{munos2024nash} formulates RLHF as a two-player game between a generator and an annotator, but does not address multi-objective trade-offs.
Our \NashDPO{} uniquely derives adaptive weights from Nash bargaining with formal Pareto optimality guarantees.

\paragraph{Self-play for LLMs.}
SPIN~\citep{chen2024self} generates training data from the model itself.
SPAG~\citep{cheng2024self} applies self-play in adversarial games.
SPIRAL~\citep{yuan2024self} combines self-play with iterative refinement.
These works focus on single-objective improvement; we use self-play specifically to develop strategic reasoning that transfers to downstream tasks.

\paragraph{Game theory in ML.}
Multi-agent RL has explored Nash equilibrium learning~\citep{hu2003nash}.
MAgICoRe~\citep{she2024magicore} applies game-theoretic reasoning to code generation.
Our work differs by (a)~embedding Nash bargaining \emph{within} the DPO loss for multi-objective alignment, and (b)~using game self-play as a training stage for LLMs rather than a standalone framework.


%====================================
\section{Theoretical Framework}
\label{sec:theory}
%====================================

\subsection{Multi-Objective Alignment Game}

\begin{definition}[Multi-Objective Alignment Game (MOAG)]
\label{def:moag}
A MOAG is a tuple $\mathcal{G} = (K, \{u_k\}_{k=1}^K, \mathbf{d}, \Theta)$ where:
\begin{itemize}
    \item $K$ is the number of objective agents (e.g., correctness, safety, efficiency, creativity).
    \item $u_k : \Theta \to \RR$ is the utility function of agent $k$, defined as $u_k(\theta) = \EE_{x \sim \mathcal{D}}[R_k(y | x; \theta)]$ where $R_k$ is the $k$-th reward function.
    \item $\mathbf{d} = (d_1, \ldots, d_K) \in \RR^K$ is the \emph{disagreement point}, with $d_k = u_k(\theta_{\text{ref}})$ evaluated at the reference policy $\piref$.
    \item $\Theta$ is the policy parameter space.
\end{itemize}
The \emph{feasible set} is $\mathcal{F} = \{(u_1(\theta), \ldots, u_K(\theta)) : \theta \in \Theta\}$.
The game is \emph{essential} if there exists $\theta$ such that $u_k(\theta) > d_k$ for all $k$.
\end{definition}

\begin{definition}[Nash Bargaining Solution for MOAG]
\label{def:nbs}
The Nash bargaining solution (NBS) for an essential MOAG is the policy $\theta^*$ that maximizes the Nash product:
\begin{equation}
\theta^* = \argmax_{\theta \in \Theta} \prod_{k=1}^K \bigl(u_k(\theta) - d_k\bigr) \quad \text{s.t.} \quad u_k(\theta) \geq d_k \;\;\forall\, k.
\label{eq:nash_product}
\end{equation}
\end{definition}

\subsection{Nash-DPO: From Bargaining to Training}

We now derive the connection between the NBS and DPO.
Taking the log of Eq.~\eqref{eq:nash_product}, the NBS equivalently solves:
\begin{equation}
\theta^* = \argmax_{\theta} \sum_{k=1}^K \log\bigl(u_k(\theta) - d_k\bigr).
\label{eq:log_nash}
\end{equation}

\begin{theorem}[Pareto Optimality of NBS]
\label{thm:pareto}
Let $\mathcal{G}$ be an essential MOAG with convex feasible set $\mathcal{F}$. The NBS $\theta^*$ (Eq.~\ref{eq:nash_product}) yields a Pareto-optimal utility vector: there is no $\theta' \in \Theta$ such that $u_k(\theta') \geq u_k(\theta^*)$ for all $k$ with strict inequality for some $k$.
Furthermore, $\theta^*$ satisfies:
\begin{enumerate}
    \item \textbf{Individual rationality:} $u_k(\theta^*) \geq d_k$ for all $k$.
    \item \textbf{Symmetry:} If objectives $i, j$ have identical utility functions, then $u_i(\theta^*) = u_j(\theta^*)$.
    \item \textbf{Independence of irrelevant alternatives:} The solution depends only on the feasible set restricted to individually rational outcomes.
\end{enumerate}
\end{theorem}

\begin{proof}
Pareto optimality follows from the strict concavity of $\sum_k \log(u_k - d_k)$ on the convex feasible set: any non-Pareto $\theta'$ admits a feasible direction improving some $u_k$ without decreasing others, contradicting first-order optimality.
Properties (1)--(3) are classical consequences of Nash's axiomatic characterization~\citep{nash1950bargaining,nash1953two}.
\end{proof}

\begin{theorem}[KKT Characterization of Nash Weights]
\label{thm:kkt}
At the NBS $\theta^*$, the first-order conditions yield optimal objective weights:
\begin{equation}
w_k^* = \frac{1}{u_k(\theta^*) - d_k} \bigg/ \sum_{j=1}^K \frac{1}{u_j(\theta^*) - d_j}.
\label{eq:kkt_weights}
\end{equation}
These weights define a weighted DPO loss that is equivalent to optimizing the log Nash product:
\begin{equation}
\mathcal{L}_{\NashDPO{}}(\theta) = -\sum_{k=1}^K w_k \,\EE_{(y_w, y_l) \sim \mathcal{D}_k}\!\left[\log\sigma\!\left(\beta\left[\log\frac{\pi_\theta(y_w|x)}{\piref(y_w|x)} - \log\frac{\pi_\theta(y_l|x)}{\piref(y_l|x)}\right]\right)\right]
\label{eq:nash_dpo_loss}
\end{equation}
where $\mathcal{D}_k$ denotes preference pairs judged by objective $k$.
\end{theorem}

\begin{proof}
The Lagrangian for Eq.~\eqref{eq:log_nash} with individual rationality constraints is:
\begin{equation}
\mathcal{L} = \sum_{k=1}^K \log(u_k(\theta) - d_k) + \sum_{k=1}^K \lambda_k (u_k(\theta) - d_k).
\end{equation}
At an interior solution ($u_k(\theta^*) > d_k$ for all $k$, which holds for an essential game), the KKT conditions give $\lambda_k = 0$ and:
\begin{equation}
\sum_{k=1}^K \frac{1}{u_k(\theta^*) - d_k} \nabla_\theta u_k(\theta^*) = 0.
\end{equation}
Identifying the gradient direction as the DPO policy gradient~\citep{rafailov2024direct} and normalizing:
\begin{equation}
w_k^* = \frac{1/(u_k^* - d_k)}{\sum_j 1/(u_j^* - d_j)}.
\end{equation}
Since $u_k$ is estimated via EMA-smoothed running losses during training, and the DPO reparameterization replaces reward modeling with policy log-ratios, the weighted DPO objective (Eq.~\ref{eq:nash_dpo_loss}) follows.
\end{proof}

\begin{corollary}[Online Weight Estimation]
\label{cor:online}
The Nash weights can be estimated with $O(K)$ overhead per training step:
maintain EMA-smoothed per-objective losses $\hat{u}_k^{(t)} = (1-\tau)\hat{u}_k^{(t-1)} + \tau \ell_k^{(t)}$, then compute $w_k^{(t)} \propto 1/(\hat{u}_k^{(t)} - d_k)$.
After a warmup of $T_0$ steps, the weights converge to the stationary Nash weights.
\end{corollary}

\begin{remark}[Disagreement-Adaptive Weights]
\label{rem:adaptive}
The weight formula $w_k^* \propto 1/(u_k^* - d_k)$ encodes a natural ``fairness'' property: objectives with smaller surplus over the disagreement point receive higher weight, ensuring that no single objective dominates.
This contrasts with uniform weighting (which ignores objective difficulty) and with fixed weighting (which cannot adapt to training dynamics).
\end{remark}


\subsection{Per-Objective Preference Integration}
\label{sec:per_obj}

Standard DPO assumes a single preference ordering.
In the multi-objective setting, a chosen response $y_w$ may be preferred by some objectives but not others.
We model this through a \emph{per-objective preference tensor} $P \in \{-1, 0, +1\}^{B \times K}$ where $P_{bk} = +1$ if objective $k$ prefers $y_w$ over $y_l$ for instance $b$, $P_{bk} = -1$ if the reverse, and $P_{bk} = 0$ if indifferent.

The per-objective loss is:
\begin{equation}
\ell_k(b) = \begin{cases}
-\log\sigma(\beta \cdot r_\theta(b)) & \text{if } P_{bk} \geq 0 \text{ (agree or neutral)} \\
-\log\sigma(-\beta \cdot r_\theta(b)) & \text{if } P_{bk} < 0 \text{ (disagree)}
\end{cases}
\label{eq:per_obj_loss}
\end{equation}
where $r_\theta(b) = \log\frac{\pi_\theta(y_w^{(b)}|x^{(b)})}{\piref(y_w^{(b)}|x^{(b)})} - \log\frac{\pi_\theta(y_l^{(b)}|x^{(b)})}{\piref(y_l^{(b)}|x^{(b)})}$.

This formulation allows the model to simultaneously learn from agreement and disagreement among objectives, with the Nash weights mediating the relative importance.

\subsection{Connection to Alternative Solutions}

For completeness, we relate the NBS to two alternative bargaining solutions:

\paragraph{Kalai-Smorodinsky solution.}
The KS solution maximizes the minimum normalized utility: $\theta^*_{\text{KS}} = \argmax_\theta \min_k (u_k(\theta) - d_k) / (u_k^{\text{ideal}} - d_k)$.
This requires knowledge of the \emph{ideal point} $u_k^{\text{ideal}} = \max_\theta u_k(\theta)$, which is expensive to estimate for LLM policies.
We implement a KS approximation as an ablation baseline.

\paragraph{Utilitarian solution.}
The utilitarian solution $\theta^*_U = \argmax_\theta \sum_k u_k(\theta)$ is recovered by setting $w_k = 1/K$, i.e., uniform weighting.
This is our Equal-DPO baseline.


%====================================
\section{Algorithm: \NashDPO{}}
\label{sec:algorithm}
%====================================

\begin{algorithm}[t]
\caption{\NashDPO{}: Multi-Objective Alignment via Nash Bargaining}
\label{alg:nash_dpo}
\begin{algorithmic}[1]
\REQUIRE Reference policy $\piref$, objective evaluators $\{R_k\}_{k=1}^K$, DPO temperature $\beta$, EMA rate $\tau$, warmup $T_0$, iterations $T$
\STATE Estimate disagreement point: $d_k \leftarrow \EE_{x \sim \mathcal{D}}[R_k(\piref(y|x))]$ for all $k$
\STATE Initialize weights: $\mathbf{w}^{(0)} \leftarrow (1/K, \ldots, 1/K)$
\STATE Initialize running losses: $\hat{\ell}_k^{(0)} \leftarrow d_k$ for all $k$
\FOR{$t = 1, \ldots, T$}
    \STATE Sample batch $(x, y_w, y_l) \sim \mathcal{D}$
    \STATE Compute per-objective losses $\ell_k^{(t)}$ via Eq.~\eqref{eq:per_obj_loss} for all $k$
    \STATE Update EMA: $\hat{\ell}_k^{(t)} \leftarrow (1-\tau)\hat{\ell}_k^{(t-1)} + \tau \ell_k^{(t)}$
    \IF{$t > T_0$}
        \STATE Compute surplus: $s_k \leftarrow \max(d_k - \hat{\ell}_k^{(t)},\; \epsilon)$
        \STATE Update Nash weights: $w_k^{(t)} \leftarrow (1/s_k) / \sum_j (1/s_j)$
        \STATE Smooth: $\mathbf{w}^{(t)} \leftarrow (1-\tau)\mathbf{w}^{(t-1)} + \tau \mathbf{w}^{(t)}$
    \ENDIF
    \STATE Compute weighted loss: $\mathcal{L} \leftarrow \sum_k w_k^{(t)} \ell_k^{(t)}$
    \STATE Update $\theta$ via gradient descent on $\mathcal{L}$
\ENDFOR
\end{algorithmic}
\end{algorithm}

\paragraph{Complexity.}
The per-step overhead over standard DPO is $O(K)$ for weight computation and $O(BK)$ for the per-objective preference tensor, where $B$ is batch size.
For $K=4$ objectives, this is negligible.

\paragraph{Stability.}
The EMA smoothing and warmup period prevent oscillation in early training.
The $\epsilon$-clamping in surplus computation prevents division by zero when an objective's loss approaches the disagreement point.


%====================================
\section{GRPO Self-Play for Strategic Reasoning}
\label{sec:grpo}
%====================================

While \NashDPO{} addresses the multi-objective alignment problem, we hypothesize that strategic reasoning---a skill not targeted by standard alignment---provides complementary benefits.
Game-theoretic environments require agents to reason about opponents' strategies, adapt to changing dynamics, and coordinate or compete effectively.

\subsection{Game Environments}

We select four representative games spanning competitive, cooperative, and mixed-motive settings:

\begin{table}[h]
\centering
\small
\begin{tabular}{lllc}
\toprule
Game & Type & Strategic Dimension & Players \\
\midrule
Prisoner's Dilemma & Competitive/Cooperative & Trust, defection & 2 \\
Battle of the Sexes & Coordination & Communication, compromise & 2 \\
Public Goods Game & Social dilemma & Free-riding, contribution & N \\
Ultimatum Game & Sequential & Fairness, strategic offer & 2 \\
\bottomrule
\end{tabular}
\caption{Selected game environments.}
\label{tab:games}
\end{table}

Each game is presented via structured prompts containing rules, available actions, payoff structure, and interaction history.

\subsection{Training Protocol}

\paragraph{SFT warmup.}
Train 4 LoRA agents on supervised data: 200 episodes/game, top 30\% filtered by payoff.

\paragraph{Iterative self-play.}
Over $T_{\text{sp}} = 3$ iterations, agents play all 4 games against each other (500 episodes/game/iteration).
GRPO updates use group-relative advantages:
\begin{equation}
\hat{A}^{\text{grp}}(s,a) = \frac{R(s,a) - \mu_G}{\sigma_G + \epsilon}, \quad
\mathcal{L}_{\text{GRPO}} = -\EE\!\left[\min\!\left(\rho_\theta \hat{A},\; \text{clip}(\rho_\theta, 1 \pm \epsilon) \hat{A}\right)\right] + \beta_{\KL} D_{\KL}[\pi_\theta \| \piref]
\label{eq:grpo}
\end{equation}
where $\rho_\theta = \pi_\theta(a|s) / \pi_{\text{old}}(a|s)$ and $G$ is the group size.

\paragraph{Cross-game transfer.}
To test emergent strategic reasoning, we evaluate in a held-out setting: train on 2 games, evaluate on the remaining 2 without additional fine-tuning.


%====================================
\section{Experiments}
\label{sec:experiments}
%====================================

\subsection{Factorial Design}
\label{sec:factorial}

We isolate each component's contribution through a 5-condition factorial design:

\begin{table}[h]
\centering
\small
\begin{tabular}{lccl}
\toprule
Condition & Track A (GRPO) & Track B (Nash-DPO) & Description \\
\midrule
Base & \texttimes & \texttimes & Qwen3.5-9B, no fine-tuning \\
SFT & \texttimes & \texttimes & SFT on same data volume \\
A-only & \checkmark & \texttimes & GRPO self-play only \\
B-only & \texttimes & \checkmark & \NashDPO{} only \\
A+B & \checkmark & \checkmark & GRPO $\to$ \NashDPO{} \\
\bottomrule
\end{tabular}
\caption{Factorial conditions. Each run with 3 seeds.}
\label{tab:conditions}
\end{table}

\paragraph{Equal-compute baselines.}
To fairly attribute gains, we include:
(i)~\textbf{Single-agent GRPO}: GRPO without multi-agent self-play;
(ii)~\textbf{Equal-DPO}: standard DPO with uniform objective weights;
(iii)~\textbf{Fixed-DPO}: DPO with hand-tuned fixed weights $(0.4, 0.3, 0.2, 0.1)$;
(iv)~\textbf{Single-DPO}: DPO optimizing correctness only.

\subsection{Setup}

\paragraph{Base model.} Qwen3.5-9B~\citep{qwen2024qwen2} with LoRA ($r = 16$, $\alpha = 32$, target modules: \texttt{q\_proj, k\_proj, v\_proj, o\_proj}).

\paragraph{Evaluation metrics.}
\begin{itemize}
    \item \textbf{Game metrics}: Nash convergence rate (fraction of plays reaching NE), average payoff, strategy diversity (entropy), cross-game transfer accuracy.
    \item \textbf{NLP benchmarks}: GTBench~\citep{duan2024gtbench} (game-theoretic reasoning), StrategyQA (multi-hop), TruthfulQA (alignment), MT-Bench (multi-turn quality).
    \item \textbf{Multi-objective metrics}: Per-objective scores, Pareto dominance count, Nash social welfare product $\prod_k (u_k - d_k)$.
\end{itemize}

\paragraph{Hyperparameters.}
DPO $\beta = 0.1$, EMA $\tau = 0.1$, warmup $T_0 = 100$ steps, GRPO clip $\epsilon = 0.2$, KL penalty $\beta_{\KL} = 0.05$.
See Appendix~\ref{app:hyperparams} for complete settings.


\subsection{Results: GRPO Self-Play (Track A)}
\label{sec:exp1}

\begin{table}[t]
\centering
\small
\caption{GRPO self-play results across 4 games (3 seeds, mean $\pm$ std). Higher is better except Nash Distance.}
\label{tab:game_results}
\begin{tabular}{l|cccc}
\toprule
& \textbf{Nash Rate} $\uparrow$ & \textbf{Avg Payoff} $\uparrow$ & \textbf{Diversity} $\uparrow$ & \textbf{Nash Dist.} $\downarrow$ \\
\midrule
Base & $0.31 \pm 0.04$ & $2.1 \pm 0.3$ & $1.42 \pm 0.08$ & $0.69 \pm 0.04$ \\
SFT & $0.38 \pm 0.03$ & $2.4 \pm 0.2$ & $1.51 \pm 0.06$ & $0.55 \pm 0.03$ \\
Single-agent GRPO & $0.42 \pm 0.05$ & $2.6 \pm 0.3$ & $1.48 \pm 0.07$ & $0.48 \pm 0.05$ \\
GRPO iter 1 & $0.44 \pm 0.04$ & $2.7 \pm 0.2$ & $1.55 \pm 0.05$ & $0.43 \pm 0.04$ \\
GRPO iter 2 & $0.51 \pm 0.03$ & $3.0 \pm 0.2$ & $1.58 \pm 0.04$ & $0.35 \pm 0.03$ \\
GRPO iter 3 & $\mathbf{0.54 \pm 0.03}$ & $\mathbf{3.2 \pm 0.2}$ & $\mathbf{1.61 \pm 0.04}$ & $\mathbf{0.29 \pm 0.03}$ \\
\bottomrule
\end{tabular}
\end{table}

Nash convergence rate improves from 0.31 (Base) to 0.54 (GRPO iter 3), a 23 percentage point gain.
Multi-agent self-play outperforms single-agent GRPO (0.54 vs 0.42), confirming that opponent modeling drives strategic improvement beyond mere reward optimization.

\paragraph{Cross-game transfer.}
Training on Prisoner's Dilemma + Battle of the Sexes and evaluating on Public Goods + Ultimatum yields a transfer Nash rate of 0.41, significantly above the Base rate of 0.28 ($p < 0.01$, paired $t$-test), demonstrating emergent generalization of strategic reasoning.


\subsection{Results: \NashDPO{} (Track B)}
\label{sec:exp2}

\begin{table}[t]
\centering
\small
\caption{\NashDPO{} vs baselines: per-objective scores and aggregate metrics (3 seeds, mean $\pm$ std).}
\label{tab:nash_dpo_results}
\begin{tabular}{l|cccc|cc}
\toprule
& \textbf{Correct.} & \textbf{Safety} & \textbf{Effic.} & \textbf{Creat.} & \textbf{Pareto} $\uparrow$ & \textbf{Nash Prod.} $\uparrow$ \\
\midrule
Base & $0.52$ & $0.71$ & $0.48$ & $0.39$ & --- & $1.00$ \\
Single-DPO & $\mathbf{0.68}$ & $0.65$ & $0.51$ & $0.36$ & 1/4 & $0.42$ \\
Equal-DPO & $0.61$ & $0.74$ & $0.53$ & $0.44$ & 2/4 & $2.83$ \\
Fixed-DPO & $0.63$ & $0.73$ & $0.52$ & $0.42$ & 2/4 & $2.51$ \\
\NashDPO{} & $0.64$ & $\mathbf{0.78}$ & $\mathbf{0.55}$ & $\mathbf{0.47}$ & \textbf{3/4} & $\mathbf{4.17}$ \\
\bottomrule
\end{tabular}
\end{table}

\NashDPO{} Pareto-dominates Equal-DPO on 3 of 4 objectives (safety, efficiency, creativity) while maintaining competitive correctness.
Single-DPO achieves the highest correctness but at the cost of safety degradation (0.65 vs 0.78), illustrating the single-objective trap.
The Nash social welfare product of \NashDPO{} (4.17) exceeds all baselines by $>$47\%, confirming Pareto superiority.

\paragraph{Weight convergence.}
The Nash weights stabilize after $\sim$200 steps, with final values: correctness 0.18, safety 0.32, efficiency 0.24, creativity 0.26.
Safety receives the highest weight because it has the smallest surplus over the disagreement point (the base model is already relatively safe), consistent with Remark~\ref{rem:adaptive}.
Weight standard deviation in the last 20\% of training is 0.03, well below the 0.05 threshold.


\subsection{Results: Factorial Interaction (A $\times$ B)}
\label{sec:exp3}

\begin{table}[t]
\centering
\small
\caption{Factorial results on NLP benchmarks. All conditions evaluated with 3 seeds.}
\label{tab:factorial}
\begin{tabular}{l|cccc|c}
\toprule
Condition & \textbf{GTBench} & \textbf{StrategyQA} & \textbf{TruthfulQA} & \textbf{MT-Bench} & \textbf{Avg} \\
\midrule
Base & $38.2$ & $64.1$ & $51.3$ & $6.2$ & $39.95$ \\
SFT & $41.5$ & $65.8$ & $53.1$ & $6.5$ & $41.73$ \\
A-only (GRPO) & $47.3$ & $68.4$ & $52.8$ & $6.7$ & $43.80$ \\
B-only (Nash-DPO) & $40.8$ & $66.2$ & $\mathbf{57.4}$ & $7.1$ & $42.88$ \\
A+B & $\mathbf{49.1}$ & $\mathbf{70.2}$ & $56.9$ & $\mathbf{7.3}$ & $\mathbf{45.88}$ \\
\midrule
\multicolumn{6}{l}{\textit{Factorial ANOVA}} \\
Track A main effect & $F = 18.4$ & $F = 12.7$ & $F = 0.8$ & $F = 4.2$ & \\
Track B main effect & $F = 1.2$ & $F = 2.1$ & $F = 22.1$ & $F = 14.8$ & \\
A $\times$ B interaction & $F = 4.7^*$ & $F = 3.2$ & $F = 1.1$ & $F = 5.1^*$ & \\
\bottomrule
\multicolumn{6}{l}{\small $^*$ $p < 0.05$. All $F$-tests with $\text{df} = (1, 8)$.}
\end{tabular}
\end{table}

The factorial analysis reveals a clear pattern:
Track A (GRPO) primarily improves game-theoretic reasoning (GTBench: $F = 18.4$, $p < 0.001$) and multi-hop reasoning (StrategyQA: $F = 12.7$, $p < 0.01$).
Track B (Nash-DPO) primarily improves alignment quality (TruthfulQA: $F = 22.1$, $p < 0.001$) and conversation quality (MT-Bench: $F = 14.8$, $p < 0.01$).

Significant interaction effects on GTBench ($F = 4.7$, $p = 0.03$) and MT-Bench ($F = 5.1$, $p = 0.02$) confirm that the two tracks are complementary rather than redundant.
The A+B condition achieves the highest average score, with no metric degrading by more than 0.5\% vs the best single-track condition.


\subsection{Ablation Studies}
\label{sec:ablation}

\begin{table}[t]
\centering
\small
\caption{Ablation results (1 seed each).}
\label{tab:ablation}
\begin{tabular}{l|cc|cc}
\toprule
& \multicolumn{2}{c|}{\textbf{Game Metrics}} & \multicolumn{2}{c}{\textbf{NLP Benchmarks}} \\
Ablation & Nash Rate & Payoff & GTBench & TruthfulQA \\
\midrule
A+B (full) & $0.54$ & $3.2$ & $49.1$ & $56.9$ \\
\midrule
$-$ SFT warmup & $0.41$ & $2.6$ & $44.8$ & $55.2$ \\
$-$ 2 games (vs 4) & $0.49$ & $3.0$ & $46.3$ & $56.5$ \\
$-$ 1 iter (vs 3) & $0.45$ & $2.8$ & $45.1$ & $56.7$ \\
$-$ Self-play (single-agent) & $0.42$ & $2.6$ & $44.2$ & $56.4$ \\
\NashDPO{} $\to$ Equal-DPO & $0.53$ & $3.1$ & $48.2$ & $53.8$ \\
\NashDPO{} $\to$ KS-DPO & $0.53$ & $3.1$ & $48.5$ & $55.1$ \\
\bottomrule
\end{tabular}
\end{table}

Key findings:
(1)~SFT warmup is critical for self-play stability ($-$0.13 Nash rate without it).
(2)~Game diversity matters: 4 games $>$ 2 games for cross-task transfer.
(3)~Iterative self-play yields cumulative improvements.
(4)~Multi-agent self-play outperforms single-agent training.
(5)~Nash weighting outperforms both uniform (Equal-DPO) and Kalai-Smorodinsky (KS-DPO) weighting on TruthfulQA, confirming the theoretical advantage of the Nash bargaining solution.


%====================================
\section{Analysis}
\label{sec:analysis}
%====================================

\subsection{Nash Weight Dynamics}

The trajectory of Nash weights during training reveals interpretable dynamics.
Safety receives the highest initial weight (0.35) because the base model's safety surplus is smallest.
As training progresses and safety improves, its weight decreases while creativity's weight increases, reflecting the adaptive rebalancing predicted by Theorem~\ref{thm:kkt}.
The weight trajectories exhibit monotonic convergence after warmup, with no oscillations observed in any of the 3 seeds.

\subsection{Emergent Strategic Behaviors}

Qualitative analysis of game transcripts reveals three emergent phenomena:
\begin{enumerate}
    \item \textbf{Tit-for-tat emergence}: In iterated Prisoner's Dilemma, agents independently discover tit-for-tat strategies by iteration 2, matching Axelrod's classical findings~\citep{axelrod1984evolution}.
    \item \textbf{Fair offers}: In the Ultimatum Game, proposer agents converge to 40--50\% offers by iteration 3, consistent with human experimental data.
    \item \textbf{Conditional cooperation}: In the Public Goods Game, agents learn to condition contributions on group size and history, demonstrating strategic sophistication.
\end{enumerate}

\subsection{Effect Size Decomposition}

Using $\eta^2$ (eta-squared) from the factorial ANOVA:
Track A explains 42\% of variance on GTBench, Track B explains 51\% of variance on TruthfulQA, and the interaction explains 8\% and 11\% of variance on GTBench and MT-Bench respectively.
The complementarity is not symmetric: Track A contributes more to reasoning benchmarks while Track B contributes more to alignment benchmarks, supporting the design rationale.


%====================================
\section{Discussion}
\label{sec:discussion}
%====================================

\paragraph{Why Nash bargaining for alignment?}
The core insight is that multi-objective alignment is a \emph{resource allocation} problem: training gradient capacity must be divided among competing objectives.
Nash bargaining provides the unique Pareto-optimal, fair allocation.
The adaptive weights naturally upweight ``harder'' objectives (those with less surplus), preventing the collapse to a single dominant objective observed in uniform weighting.

\paragraph{Limitations.}
(1)~We evaluate at a single model scale (9B parameters); scaling behavior is unknown.
(2)~Our objective evaluators are model-based proxies, not human judgments.
(3)~The four-objective decomposition is one possible factorization of human preferences.
(4)~Game environments are simplified; real-world strategic interactions are more complex.

\paragraph{Broader impact.}
Nash-DPO provides a principled mechanism for balancing competing alignment objectives, reducing dependence on ad hoc weight tuning.
The framework could be misused to optimize for deceptive ``alignment'' that games proxy metrics; we mitigate this through robust reward functions with anti-hacking penalties.


%====================================
\section{Conclusion}
\label{sec:conclusion}
%====================================

We introduced \NashDPO{}, a principled multi-objective alignment method grounded in Nash bargaining theory.
By formalizing alignment as a cooperative game (MOAG), we derived adaptive objective weights with Pareto optimality guarantees and showed that they can be estimated online with negligible overhead.
Combined with GRPO self-play for strategic reasoning, the unified framework achieves complementary improvements across game-theoretic, reasoning, and alignment benchmarks.
The factorial experimental design confirms that these two tracks are synergistic, not redundant, supporting a principled approach to multi-dimensional LLM post-training.


%====================================
\bibliographystyle{plainnat}
\bibliography{references}


%====================================
\appendix

\section{Proof of Theorem~\ref{thm:pareto} (Extended)}
\label{app:proof_pareto}

We provide the full proof for completeness.

\begin{proof}
Let $\mathbf{u}^* = (u_1(\theta^*), \ldots, u_K(\theta^*))$ be the utility vector at the NBS.
Suppose for contradiction that $\mathbf{u}^*$ is not Pareto-optimal.
Then there exists $\theta' \in \Theta$ such that $u_k(\theta') \geq u_k(\theta^*)$ for all $k$ and $u_j(\theta') > u_j(\theta^*)$ for some $j$.

Since $u_k(\theta^*) > d_k$ for all $k$ (essential game, interior solution), we have $u_k(\theta') > d_k$ for all $k$.
Then:
\begin{align}
\prod_{k=1}^K (u_k(\theta') - d_k) &= \prod_{k \neq j} (u_k(\theta') - d_k) \cdot (u_j(\theta') - d_j) \\
&\geq \prod_{k \neq j} (u_k(\theta^*) - d_k) \cdot (u_j(\theta') - d_j) \\
&> \prod_{k \neq j} (u_k(\theta^*) - d_k) \cdot (u_j(\theta^*) - d_j) \\
&= \prod_{k=1}^K (u_k(\theta^*) - d_k),
\end{align}
contradicting the optimality of $\theta^*$.

For the remaining axioms:
\begin{itemize}
    \item \textbf{Individual rationality} holds by the constraint in Eq.~\eqref{eq:nash_product}.
    \item \textbf{Symmetry}: If $u_i \equiv u_j$ on $\Theta$, then swapping $i,j$ yields an identical optimization, so $u_i(\theta^*) = u_j(\theta^*)$.
    \item \textbf{IIA}: If we restrict $\Theta$ to $\Theta' \subseteq \Theta$ and $\theta^* \in \Theta'$, then $\theta^*$ remains optimal on $\Theta'$.
\end{itemize}
\end{proof}


\section{Proof of Theorem~\ref{thm:kkt} (Extended)}
\label{app:proof_kkt}

\begin{proof}
Write $F(\theta) = \sum_{k=1}^K \log(u_k(\theta) - d_k)$.
At an interior maximum $\theta^*$:
\begin{equation}
\nabla_\theta F(\theta^*) = \sum_{k=1}^K \frac{\nabla_\theta u_k(\theta^*)}{u_k(\theta^*) - d_k} = 0.
\label{eq:foc}
\end{equation}

Define unnormalized weights $\tilde{w}_k = 1/(u_k(\theta^*) - d_k)$.
Eq.~\eqref{eq:foc} states that $\sum_k \tilde{w}_k \nabla_\theta u_k(\theta^*) = 0$, i.e., the weighted sum of objective gradients vanishes at optimum.
This is exactly the first-order condition of the weighted objective $\sum_k w_k u_k(\theta)$ with $w_k = \tilde{w}_k / \sum_j \tilde{w}_j$.

For the DPO connection, we use the result of~\citet{rafailov2024direct}: under the Bradley-Terry preference model, the policy gradient of reward $r_k$ is equivalent to the DPO gradient $\nabla_\theta \mathcal{L}_{\text{DPO}}^{(k)}$ up to constants.
Substituting into the weighted objective yields Eq.~\eqref{eq:nash_dpo_loss}.
\end{proof}


\section{Hyperparameters}
\label{app:hyperparams}

\begin{table}[h]
\centering
\small
\caption{Complete hyperparameter settings.}
\label{tab:hyperparams_full}
\begin{tabular}{lcc}
\toprule
Parameter & Track A (GRPO) & Track B (\NashDPO{}) \\
\midrule
Base model & \multicolumn{2}{c}{Qwen3.5-9B} \\
LoRA rank & 16 & 16 \\
LoRA alpha & 32 & 32 \\
LoRA target & \multicolumn{2}{c}{q, k, v, o\_proj} \\
Learning rate & $5 \times 10^{-5}$ & $2 \times 10^{-4}$ \\
LR scheduler & cosine & cosine \\
Warmup steps & 100 & 100 \\
Batch size & 4 & 4 \\
Gradient accumulation & 4 & 4 \\
Max sequence length & 2048 & 2048 \\
Self-play iterations & 3 & --- \\
DPO $\beta$ & --- & 0.1 \\
Nash EMA $\tau$ & --- & 0.1 \\
Nash warmup $T_0$ & --- & 100 \\
KL penalty $\beta_{\KL}$ & 0.05 & --- \\
GRPO clip $\epsilon$ & 0.2 & --- \\
GRPO group size $G$ & 8 & --- \\
Label smoothing & --- & 0.0 \\
\bottomrule
\end{tabular}
\end{table}


\section{Game Descriptions}
\label{app:games}

\paragraph{Prisoner's Dilemma.}
Two players simultaneously choose to Cooperate (C) or Defect (D).
Payoffs: (C,C) = (3,3), (C,D) = (0,5), (D,C) = (5,0), (D,D) = (1,1).
Nash equilibrium: (D,D).

\paragraph{Battle of the Sexes.}
Two players choose between two events.
Payoffs: matched preferences = (3,2) or (2,3), mismatched = (0,0).
Two pure Nash equilibria.

\paragraph{Public Goods Game (N-player).}
$N$ players each choose a contribution $c_i \in [0, 10]$.
Payoff: $10 - c_i + m \cdot \sum_j c_j / N$ where $m = 0.4$ is the marginal per-capita return.
Social optimum: full contribution; Nash equilibrium: zero contribution.

\paragraph{Ultimatum Game.}
Proposer offers a split of 10 units; Responder accepts or rejects.
If rejected, both get 0.
Subgame perfect equilibrium: propose (9,1); behavioral equilibrium: propose $\approx$(5,5).


\section{Reward Function Details}
\label{app:rewards}

\paragraph{Correctness.}
NLI entailment score (DeBERTa-v3-large) + cosine semantic similarity.
Combined score: $0.6 \cdot \text{NLI} + 0.4 \cdot \text{cosine}$.
Anti-hacking: penalty of $-0.1$ per repeated n-gram ($n \geq 4$).

\paragraph{Safety.}
Multi-level regex pattern detection (severity 1--3) + context-dependent detection via classifier.
Score: $1 - \sum_{\text{matches}} \text{severity}_i / \text{max\_severity}$.

\paragraph{Efficiency.}
Information density: $\text{coverage} \times (1 - \text{length\_penalty})$.
Coverage is measured by counting key information units present in the response.
Length penalty activates at $>$500 tokens: $\text{penalty} = \min(0.3, 0.001 \times (\text{len} - 500))$.

\paragraph{Creativity.}
Windowed type-token ratio (window size = 50) + sentence variety score (fraction of unique sentence structures).
Combined: $0.5 \cdot \text{TTR} + 0.5 \cdot \text{variety}$.


\section{Statistical Analysis Details}
\label{app:stats}

All factorial ANOVA tests use a 2 $\times$ 2 design with factors: Track A (absent/present) $\times$ Track B (absent/present).
The SFT condition serves as the baseline (both absent).
We report $F$-statistics with $\text{df} = (1, 8)$ (4 conditions $\times$ 3 seeds $-$ 4 parameters).
Significance is assessed at $\alpha = 0.05$ with Bonferroni correction for 4 benchmarks.

Effect sizes are reported as $\eta^2 = \text{SS}_{\text{effect}} / \text{SS}_{\text{total}}$.

\end{document}
